% Round 31 — Chương 4 phần A v4 và phần B v7 là nguồn các kết quả tổng kết.
% Các mã EV thuộc sổ chính và các phụ lục bằng chứng 1–4.
\chapter{Kết luận}
\label{chuong:ketluan}

Khóa luận đã đặc tả, cài đặt và đo trần phần dời lời hứa giờ trả do điều phối của từng hành khách. Chương này tổng kết kết quả, hạn chế và hướng phát triển.
% EV: A01; A06; B09; X3-12

\section{Tổng kết theo ba phần}
\label{muc:c5-tongket}

Với phần tính cho điều phối, luật C giới hạn tổng độ lớn các lần dời giờ trả của mỗi khách qua cả chuyến. Mốc so sánh là giờ trả tính lại khi giữ nguyên kế hoạch cũ, tiếp nối mốc trong công thức 7 của Hörl và Zwick~\cite{horl2022}. Trên 34 kịch bản với 3.672 yêu cầu, luật C trần 30 giây phục vụ 2.338 lượt, so với 2.434 khi không trần: ít hơn 96 lượt. Không khách nào trong 2.338 khách vượt trần; tổng dời lớn nhất là 29,2 giây. So với lời hứa đầu không dự báo, số lượt không trễ quá một phút là 2.262 thay vì 2.231 trên 3.672 yêu cầu. Luật hạn chót 30 giây phục vụ cùng số lượt, nhưng có khách bị phá lời hứa của chính luật ở 28 trong 34 kịch bản. Nới trần lên 120 giây cho 2.387 lượt phục vụ, 2.255 lượt không trễ quá một phút và 2.366 lượt không trễ quá hai phút trên 3.672 yêu cầu. Đây là một điểm đo thêm trên đường đánh đổi. Vế cộng dồn chỉ riêng nó giải thích 3 trong 524 bản ghi chặn ở mức này; tác dụng lên kết quả phục vụ chưa được đo.
% EV: A06; A02; J03; C01; C02; C20; C08; C13; N-21; N-06; L1a-02; K36; N-08; L1a-06; E09

Với phần được hồ sơ tắc đường lịch sử giải thích, lời hứa theo dự báo giảm số khách trễ quá một phút từ 68 xuống 2 trong 1.471 khách của 21 kịch bản ngày thường, gộp theo khách. Đây là phép tính lại lời hứa trên các lần chạy đã ghi. Trên cả 3.672 yêu cầu của tập chính, luật C có dự báo cho 2.328 lượt không trễ quá một phút. Điều phối không trần có dự báo và đệm thêm 30 giây cũng cho đúng 2.328 lượt; ở ngưỡng hai phút, cách đệm cho 2.366 lượt, nhiều hơn 2.338 của luật C. Đổi lại, lời hứa dài thêm trung bình 37,2 giây thay vì 7,4 giây, gộp theo khách được phục vụ. Số khách tới sớm quá một phút là 44 trong 2.434 với cách đệm và 4 trong 2.338 với luật C. Vì vậy, số lượt không trễ ở một ngưỡng chưa đủ để chọn giữa trần và đệm.
% EV: F04; F17; F03; C09; C10; N-16; N-18; K06; K51

Với phần còn lại, sổ theo từng khách cho biết dự báo đang dùng chưa giải thích được bao nhiêu độ trễ. Trong 6.217,4 giây trễ của 68 khách ngày thường nói trên, phần được hồ sơ lịch sử giải thích là 5.352,6 giây, phần còn lại là 783,0 giây. Ở chiều tuyết, 41 trong 601 khách trễ quá một phút so với lời hứa đầu không dự báo. Tổng 3.445,9 giây trễ của họ gồm 8,6 giây tính cho điều phối, 251,5 giây được hồ sơ lịch sử giải thích và 3.185,9 giây còn lại. Sổ là phép phân chia theo mô hình và dự báo đang dùng; phần còn lại không được hiểu là phần không thể biết trước.
% EV: F09; F04; F17; F13; N4B-05; K04; K19; X3-12

\section{Hạn chế}
\label{muc:c5-hanche}

Mục~\ref{muc:c4-thaoluan} trình bày chi tiết các hạn chế. 34 kịch bản thuộc 5 ngày phụ thuộc nhau, mỗi kịch bản chạy một lần. Mô phỏng có tải thấp và một hệ số tắc đường chung cho cả thành phố; chưa đo phản ứng của khách hay giá trị của trần với hãng xe. Đổi bộ chọn và luật mỗi lần một trần chưa được chạy trên quỹ đạo. Chưa tách được phần của bộ chọn và chưa đo giá trị tăng thêm của vế cộng dồn.
% EV: I01; I02; I04; I05; B11; B12; I07; I10; H12; H02; E09; K18; K08

Đóng góp còn giới hạn bởi phần kế thừa. Fielbaum và Alonso-Mora đã đo và kiểm soát bất ổn trước quyết định~\cite{fielbaum2020}; Hörl và Zwick đã tính lại tuyến cũ theo tắc đường~\cite{horl2022}. Schulz và Pfeiffer đã đặt hạn chót từ lời hứa đầu~\cite{schulz2026}. Chèn động và đánh đổi giữa phục vụ với ổn định cũng đã có tiền lệ. Khóa luận cài đặt, kiểm tra một nghĩa vụ cụ thể và đo cái giá của nó.
% EV: J06; J08; J01; J03; J10; J38; K30; K31

\section{Hướng phát triển}
\label{muc:c5-huongphattrien}

Thứ nhất, cần thử kết hợp hai thứ tự ưu tiên của bộ chọn với hai cấu hình không trần và trần 30 giây trên tập kiểm tra chính. Bộ chạy mới phải tái lập kết quả cũ trên các kịch bản kiểm tra trước khi so sánh. Phép thử cần báo cả lượt phục vụ, lượt trễ, lượt không trễ và khách đi chung để đo sự phụ thuộc của đánh đổi vào bộ chọn.
% EV: H12; H02; K18; I13; OUTLINE-V2 §7.2 (hướng đề xuất, chưa có kết quả)

Thứ hai, cần chạy luật mỗi lần một trần cạnh luật cộng dồn, và thêm hạn chót nới theo giờ trả khi giữ nguyên kế hoạch làm đối chứng. Hai phép so này kiểm riêng vai trò của lịch sử các lần dời và của mốc thời gian cập nhật.
% EV: E09; E08; A22; K27; J03; J05 (hướng đề xuất)

Thứ ba, cần kiểm dự báo cho ngày có sự kiện bằng dữ liệu có đúng tại thời điểm hứa và trên nhiều sự kiện độc lập. Số khách trễ phải được báo cùng số khách tới sớm và độ dài thêm của lời hứa.
% EV: B26; I11; F14; K20; K21; C10; C11 (hướng đề xuất)

Thứ tư, cần mở rộng phép đo sang tải cao hơn và tắc đường riêng trên từng đoạn đường. Mô hình phản ứng của khách, cùng chi phí và ngưỡng chấp nhận cụ thể của hãng xe, là dữ liệu cần có để đánh giá giá trị sử dụng của trần.
% EV: I02; I04; B12; I07; I10; D20 (hướng đề xuất)

Thứ năm, cần thử trần theo độ dài chuyến và sửa phép kiểm giờ đón để phù hợp với lời hứa theo dự báo. Kết quả phải đo trên quỹ đạo chạy thật của từng cấu hình, thay vì suy số khách được phục vụ từ phép sàng lọc.
% EV: D18; F07; I13; K11; J25 (hướng đề xuất)
